% Created 2018-05-04 Fri 16:08
% Intended LaTeX compiler: pdflatex
\documentclass[11pt]{article}
\usepackage[utf8]{inputenc}
\usepackage[T1]{fontenc}
\usepackage{graphicx}
\usepackage{grffile}
\usepackage{longtable}
\usepackage{wrapfig}
\usepackage{rotating}
\usepackage[normalem]{ulem}
\usepackage{amsmath}
\usepackage{textcomp}
\usepackage{amssymb}
\usepackage{capt-of}
\usepackage{hyperref}
\usepackage[margin=0.80in]{geometry}
\usepackage[backend=biber, style=ieee]{biblatex}
\usepackage{float}
\usepackage[super,negative]{nth}
\usepackage[capitalise]{cleveref}
\usepackage{pst-node,transparent,ragged2e}
\addbibresource{/home/nuk3/.spacemacs.d/org-files/bibliography.bib}
\DeclareFieldFormat[inproceedings]{citetitle}{\textit{#1}}
\DeclareFieldFormat[inproceedings]{title}{\textit{#1}}
\DeclareFieldFormat[misc]{citetitle}{#1}
\DeclareFieldFormat[misc]{title}{#1}
\renewcommand*{\bibpagespunct}{%
\ifentrytype{inproceedings}
{\addspace}
{\addcomma\space}}
\AtEveryCitekey{\ifuseauthor{}{\clearname{author}}}
\AtEveryBibitem{\ifuseauthor{}{\clearname{author}}}
\SpecialCoor
\author{Nyameko Lisa}
\date{13 April 2018}
\title{IPSP055 - Essential Industrial Design Law\\\medskip
\large Assignment 01 - 751243}
\hypersetup{
 pdfauthor={Nyameko Lisa},
 pdftitle={IPSP055 - Essential Industrial Design Law},
 pdfkeywords={},
 pdfsubject={},
 pdfcreator={Emacs 27.0.50 (Org mode 9.1.6)},
 pdflang={English}}
\begin{document}

\maketitle
\addvspace{110pt}
\centering{
\pnode(0.5\textwidth,-0.5\textheight){thisCenter}
\rput(thisCenter){%\transparent{0.25}
\includegraphics[width=2.7in]{/home/nuk3/course/llb/wipo-unisa/UNISACoatofArms.eps}}}

\justifying
\addvspace{110pt}
\section*{Declaration}
\label{sec:org8fb61e9}
I know that plagiarism is to use someone else’s work and pass it off as my own.\\
I know that plagiarism is wrong.\\
I confirm that this assignment is my own work.\\
I have acknowledged in the bibliography accompanying the assignment all the sources that I have used.\\
I have not directly copied without acknowledgement anything from the Internet or from any other source.\\
I have indicated every quotation and citation in a footnote or bracket linked to that quotation.\\
I have not allowed anyone else to copy my work and to pass it off as their own work.\\
I understand that if any unacknowledged copying whatsoever appears in my assignment I will receive zero per cent for the assignment.\\
I am aware of the UNISA policy on plagiarism and understand that disciplinary proceedings can be instituted against me by UNISA if I contravene this policy.\\
I indicate my understanding and acceptance of this declaration by
entering my name here under:
\begin{itemize}
\item Name: \textbf{Nyameko Lisa} (Student Number: \textbf{7874-909-3})
\end{itemize}

\subsection*{NOTE}
\label{sec:org0782d60}
Please note that footnotes will be denoted as \footnote{This is a footnote.} and will
appear at the bottom of the page.\\
References will be denoted by \cite{rsa93_designs_act} and will appear at the end of the document.
\newpage
\section{Advise whether the design will qualify for design protection and what the steps are to register it. [20]}
\label{sec:org7e2e903}

\subsection{Introduction}
\label{sec:org9449bdd}

As per the definitions of [section 2(1)]\cite{rsa93_tm_act} and given that the
Republic is both a Paris Convention as well as a World Intellectual Property
Organisation (WIPO) Member State, [article 15(1)]\cite{wto17_trips}, [article
1(2)]\cite{wipo83_paris_conve_protect_ip}, \cite{wipo25_hague_agreement}, South African statutory design
legislation is in compliance with International Conventions and legal instruments.\\

As per the provisions of \cite{rsa78_patents_act}, patent rights seek to protect
the principle underlying an inventive process or product. In contrast the South
African Designs Act No. 195 of 1993 \cite{rsa93_designs_act}, provides provisions
for the protection of only the physical form of articles of manufacture intended
to be multiplied by industrial processes, [section 14(4)]\cite{rsa93_designs_act}.

\subsection{Requirements for design protection}
\label{sec:orgf372765}

Judson Beaumont, who will henceforth be referred to as the \textbf{``proprietor''} of
the Beaver Cabinet design which will in turn be referred to as the \textbf{``article''}, as
per the provisions of [section 1(1)]\cite{rsa93_designs_act} and his Straight Line
Designs studio should be advised that the legislation makes provision for both
\uline{aesthetic designs} (Part A) as well as \uline{functional designs} (Part F)
registrations, as per the provisions of [section
14(1)(a-b)]\cite{rsa93_designs_act}.

\subsubsection{Aesthetic designs}
\label{sec:orgf221d5e}

As per the provisions of [section 1(1)]\cite{rsa93_designs_act}, which stipulates
that \textit{``...any design applied to any article, whether for the
pattern or the shape or the configuration or the ornamentation thereof, or for
any two or more of these purposes, and by whatever means is applied, having
features which appeal to and are judged solely by the eye, irrespective of the
aesthetic quality thereof...''}, as was stipulated by Judge Corbett in
\citetitle{corbett84_homecraft_steel_v_sm_hare}:
\begin{quote}
\textit{To qualify as a design, an article must have some special feature
which appeals to the eye and distinguishes it from others of its class.''}
\end{quote}
The \textbf{proprietor} would be advised to pursue aesthetic design registration for
their \textbf{article}.

\subsubsection{Functional designs}
\label{sec:orgf38ba0c}
Again as per the provisions of [section 1(1)]\cite{rsa93_designs_act}, the
definition of a function design is \textit{``...any design applied to any article,
either for the pattern or the shape or the configuration thereof, or for any two
or more of these purposes, and by whatever means it is applied, having features
which are necessitated by the function which the article to which the design is
applied, is to perform, and includes an integrated circuit topography, a mask
work and a series of mask works...''}.\\

Whist it is possible for the \textbf{proprietor} to lodge both functional and aesthetic
design registration applications in respect of their \textbf{article}, as per the
provisions of [section 14(1)(b)(ii)]\cite{rsa93_designs_act}, it is not obviously
apparent that there are any functional features of the Beaver Cabinet article
that require protection of purely functional design features or methods and
principles of it's construction [section 14(4)]\cite{rsa93_designs_act}, \emph{``that
are not commonplace in the art in question''}. Without any evidence to the
contrary, Straight Line Designs would most likely be unsuccessful in seeking
functional design registration for their \textbf{article}.

\subsubsection{Novelty requirements}
\label{sec:org6986516}

As per the provisions of [section 14(1)(a)]\cite{rsa93_designs_act}, should Judson
be able to demonstrate that his Beaver Cabinet is both \textbf{new} and \textbf{original}, he
may in the prescribed manner and on payment of the prescribed fee, apply for
aesthetic design registration. Moreover the Beaver Cabinet must be different
from or not form part of the \textbf{state of the art} within a grace period of six
months from the date of first disclosure, or release date [section
14(2)(b)]\cite{rsa93_designs_act}. Where state of the art comprises of all matter
which has been made available to the public, within or outside the Republic, by
written means, usage or any other form [section 14(3)(a)]\cite{rsa93_designs_act},
including all matter contained in pending registration applications [section
14(3)(b)]\cite{rsa93_designs_act}. Courts of the Republic have interpreted the
meaning of \emph{`original'} to coincide with that definition outlined in
\cite{rsa78_copyrightact}, and which was further elaborated on by Judge Van
Heerden in \citetitle{vheerden71_xactics_v_tailored_containers}, where he
explained that:
\begin{quote}
\textit{``I have accordingly come to the conclusion that applicant's registered design
68/346 was not new or original by reason of the fact that a design not substantially
different therefrom had already been described in a printed publication before the
date of application for registration.''}
\end{quote}

\subsection{Requirements and procedure for registration}
\label{sec:org4c873b9}
As per the provisions of [section 44]\cite{rsa93_designs_act}, should Judson have
either already registered or applied for protection of the Beaver Cabinet design
in a convention member state outside of the Republic, then as the \textbf{proprietor}
of said \textbf{article}, in accordance with \cite{rsa93_designs_act}, he may indeed
apply for registration within the Republic. Judson my also ensure that he
himself is protected by clearly articulated, binding employment agreements with
his staff of eight full-time craftspeople, to ensure that he is the unequivocal
sole proprietor of the design [section 14(7)]\cite{rsa93_designs_act}.\\

As per the provisions of [section 15]\cite{rsa93_designs_act}, there is neither
examination nor provision for opposition with regards to the Beaver Cabinet's
registrability. The Registrar will consider Judons' application and should he be
satisfied that all formal requirements have been met in the prescribed manner
[section 16]\cite{rsa93_designs_act}, he will register the Beaver Cabinet design
and issue the official Certificate of registration [section
18]\cite{rsa93_designs_act}. After the certificate is issued and the design is
published as per [section 18(b)]\cite{rsa93_designs_act}, the official file is
released to the public for inspection [section 19]\cite{rsa93_designs_act}. Judson
should expect the process to be completed within 12 months of the date of
filing. Moreover there are provisions for him to correct any clerical errors
upon application [section 26]\cite{rsa93_designs_act}.

\subsubsection{Prescribed filing requirements}
\label{sec:orgdecb084}
Judon's application for the registration of his Beaver Cabinet, must comply with
the above mentioned statutory requirements, and be accompanied by the prescribed
forms, containing all the required information including but not limited to:

\begin{itemize}
\item As the proprietor of the article's design, Judson must provide his name and
physical address,
\item Article for which design protection is to be applied (Cabinet), and the
corresponding international classification \cite{wipo79_locarno_agreement},
\item Type of Design (Aesthetic only \emph{recommended}),
\item Number, date and country of the basic application if Convention priority is to
be claimed,
\item A declaration of novelty or definitive statement outlining all the features
of the design for which protection is sought.
\end{itemize}


\printbibliography
\end{document}